%% 自動生成（phys/gen_phys.py）
\phtocgrid{\phgi{保存則と相似の地図}{4}\phgi{この本の使い方}{6}}
{\color{bkLine}\hrule height 0.35pt}\vspace{2.4mm}
\phtocchap{力-0}{クレーンは何トンまで持ち上げられる？}{荷を持ち上げるとき、クレーンのどこに力がかかる？}{}{3}{3}{8}{J}
\phtocpart{第1部}{動く：質点と質点系}{24}{1}
\phtocchap{力-1}{動きを数で書く：位置・速度・加速度}{カーブを曲がる車は、速さが同じでも加速している？}{4,5}{3}{4}{26}{C}
\phtocchap{力-2}{力の図から式へ：運動方程式を立てる}{エレベータが上へ動き出すと、体が重く感じるのはなぜ？}{3,6}{3}{4}{86}{C}
\phtocchap{力-3}{ぶつかる・噴き出す：運動量と衝突}{ロケットは、何もない宇宙でどうやって進む？}{5,8}{4}{4}{146}{C}
\phtocchap{力-4}{高さが速さに変わる：仕事とエネルギー}{ジェットコースターは、最初の山より高く上がれる？}{4,8}{4}{5}{206}{C}
\phtocpart{第2部}{回る：剛体と機構}{258}{2}
\phtocchap{力-5}{橋を支える三角形：つり合いとトラス}{鉄橋は、なぜ三角形を組み合わせてつくる？}{4,7}{3}{4}{260}{C}
\phtocchap{力-6}{回る物体：慣性モーメントと角運動量}{スケートの選手が腕をたたむと、速く回るのはなぜ？}{4,8}{4}{5}{322}{C}
\phtocchap{力-7}{転がる・ロボットの腕：平面運動と機構}{ロボットの手を、ねらった場所へ動かすには？}{4,7}{4}{5}{372}{C}
\phtocbreak
\phtocpart{第3部}{ゆがむ：材料の力学}{438}{3}
\phtocchap{力-8}{引っぱる・押す：応力とひずみ}{鉄の棒も、引っぱると少しのびる？}{2,3}{3}{4}{440}{C}
\phtocchap{力-9}{向きで変わる応力：組み合わせ応力とモールの円}{チョークをねじると、ななめに折れるのはなぜ？}{4,7}{5}{5}{494}{C}
\phtocchap{力-10}{板がたわむ：はりの曲げとたわみ}{本棚の板は、なぜまん中がいちばんたわむ？}{5,6,8}{4}{5}{540}{C}
\phtocchap{力-11}{ねじる・ふくらむ・曲がる：ねじりと座屈}{細い棒を上から押すと、急にぐにゃりと曲がるのはなぜ？}{5,6,8}{5}{5}{594}{C}
\phtocpart{第4部}{流れる：流体の力学}{642}{4}
\phtocchap{力-12}{水の中の力：圧力と浮力}{鉄でできた大きな船が、なぜ水に浮かぶ？}{5,8}{3}{4}{644}{C}
\phtocchap{力-13}{流れを数える：質量・エネルギー・運動量}{蛇口から落ちる水は、なぜ下へいくほど細くなる？}{8}{4}{5}{702}{C}
\phtocchap{力-14}{管の中の流れ：粘性・層流と乱流・圧力損失}{長い水道管の先では、なぜ水の勢いが弱くなる？}{3,6}{4}{5}{752}{C}
\phtocchap{力-15}{模型で本物を調べる：物体まわりの流れと相似則}{小さな模型の実験で、本物の飛行機にはたらく力がわかる？}{2,7,8}{5}{5}{806}{C}
\phtocchap{力-16}{コンピュータで解く力学}{車をこわさずに、ぶつかったときのようすを調べるには？}{7,8}{4}{5}{852}{R}
\par\vspace{1.5mm}
\phtocgrid{\phgi{こたえ}{882}\phgi{記号カード}{898}\phgi{物性値・定数・単位}{903}\phgi{\mbox{なぜ？}の道すじ}{906}\phgi{できるようになったことチェック表}{908}\phgi{つぎの巻へ}{909}\phgi{さくいん}{910}}
